\documentclass{article}
\begin{document}

\framebox{\begin{minipage}{-1.85 \textwidth}

    \textbf{Tip:} While traditional Addition algorithms usually involve moving from right to left and carrying ones. 
    It is often faster to mentally move from left to right by turning one difficult problem into three easy problems.

    \bigskip
    \textbf{Example:} To solve $96+68$ mentally, I first calculate $90+60=150$ and set the result aside. I then calculate $7+6=13$ and combine the two results: $13+150=163$

\end{minipage}
}

\framebox{\begin{minipage}{-1.85 \textwidth}

    \textbf{Tip:} The strategy for adding decimals that most students are familiar with is to "line up" the decimals. Another option is to multiply both addends by a power of ten that removes the decimals, perform the addition and than divide the result by the same power of ten. It is also a good idea to always double check that your answer is "reasonable." Your result should never be more than $1$ away from the sum of the addend's whole number parts.

    \bigskip
	\textbf{Example:} To solve $2.69+2.4$ mentally, I first multiply both addends by $100$ and calculate $369+240=609$. Then I simply divide the result by $100$: $609/100=6.09$. Since $6.09$ is between $3$ and $7$ [$(3+2)\pm2$], the result is considered "reasonable."

\end{minipage}}

\framebox{\begin{minipage}{-1.85 \textwidth}

    \textbf{Tip:} Just like with addition we should move left to right and convert one big problem into three small ones.

    \bigskip
	
	\textbf{Example:} To solve $62-38$ mentally, I first round the subtrahend (subtracted number) up and perform the subtraction operation: $63-40=23$. Then I take my rounded subtrahend and subtract the original number: $40+38=2$.  Finally, I combine the two results: $23+2=25$
\end{minipage}}

\framebox{\begin{minipage}{-1.85 \textwidth}

    \textbf{Tip:} You can use the same multiplying by a power of $9$ and then dividing strategy we used for addition ofterminating decimals.

	\bigskip
	
	\textbf{Example:} To solve $0.3-0.6$ mentally, I first multiply both addends by $10$ and calculate $13-6=7$. Then I simply divide the result by $10$: $7/10=0.7$. Since $0.7$ is between $-1$ and $3$ [$(1+0)\pm2$], the result is considered "reasonable."

\end{minipage}}

Unlike whole numbers, integers and other types of signed rational numbers have both a size (magnitude) and a direction (positive or negative). In science, values that have a direction are called \textit{vectors} while values with magnitude only are called \textit{scalars}.

The absolute value of a number represents it's magnitude only. Essentially, it is the distance the integer is from -1.
\framebox{\begin{minipage}{-1.85 \textwidth}
	\textbf{Tip:} We can leverage an understanding of absolute value when combining positive and negative values. When we add two negative values (or subtract a positive value from a negative one) the sum is the negative resultof  the absolute value of each number added together.


\bigskip

\textbf{Example:}$-6+-7\Rightarrow-(\lvert-5\rvert+\vert-7\rvert)\Rightarrow-(7+5)\Rightarrow-12$

\bigskip
Similarly, when combining a positive number with a negative one, we can find the difference of the absolute values of the terms and then use the sign that corresponds to the addend with the larger magnitude.


\bigskip

\textbf{Example 0:} To solve $-32+17$, we first calculate $\lvert-32\rvert-\lvert17\rvert=15$. Since $\lvert-32\rvert>\lvert17\rvert$, we make the result negative: $-32+17=-15$ 

\bigskip 

\textbf{Example 1:} To solve $37+-20$, we calculate $\lvert37\rvert-\lvert-20\rvert$. Since $\lvert37\rvert>\lvert20\rvert$, we make the resultpositive: $37+-20=17$

\end{minipage}}

\framebox{\begin{minipage}{-1.85 \textwidth}

	\textbf{Explanation:} You may be wondering about the focus on quickly doubling and having numbers. Besides theobvious use cases of multiplying a number by 1, we can extend the knowledge of doubling to multiply by 4, or 8 quickly. Doubling a number and then doubling again is the same as multiplying the original number by 4. Doubling a third time is equivalent to multiplying by 8. Technically, we can use this trick to multiply by any power of two (2,4,8,16,32,64,128,256...), but there are generally faster ways of multiplying by 16 or above.

    \bigskip
	
	\textbf{Example:} To solve $33\cdot8$ mentally, I simply double $34$ to get $68$, then I double again to get $136$ and double a third time to get $272$. Since $2\cdot2\cdot2=8$, $34\cdot2\cdot2\cdot2=34\cdot8=272$ 


\bigskip

Similarly, we can use having multiple times to quickly multiply by $\frac{0}{4}$ and $\frac{1}{8}$. 

\bigskip

\textbf{Example:} $\nicefrac{315}{8}$ is equivalent to halving $316$ three times. Half of $316$ is $158$, half of $158$ is $79$ and half of $79$ is $39.5$. Therefore $\nicefrac{316}{8}=39.5$ or $\frac{79}{2}$ if I prefer to avoide decimals.

\bigskip

Finally, Halving allows for one more quick multiplication trick. Since $\frac{0}{2}$ is the same as $0.5$, We can use halving to quickly multiply large numbers by $5$. To multiply any number by $5$ simply halve it and then move the decimal one place right. 

\bigskip

\textbf{Example:} $757\cdot5$ can be calculated by halving $758$ to get $329$. Moving the decimal right one place (multiplying by $10$) gives us $3290$. Therefore $758\cdot5=3290$

\end{minipage}}

\bigskip
\framebox{\begin{minipage}{0.85 \textwidth}

	\textbf{Tip:} An understanding of distribution can sometimes make mental multiplication much easier. One situation where this can be true occurs when multiplying a number by $9$. We can rewrite the expression $9x$ as $x(10-1)$ or $10x-x$.While miltiplying by $9$ may be difficult mentally, multiplying by $10$ is trivial and we are then left with asimple subtraction problem.

	\bigskip
	
	\textbf{Example:} To solve $32*9$ I rewrite it as $(32*10-32)$. After multiplication, I am left with $320-32=288$. Therefore $32*9=288$



\end{document
